\documentclass[12pt,spanish,mexico]{article}
\usepackage[utf8]{inputenc}
\PassOptionsToPackage{hyphens}{url}
\usepackage{hyperref}
\usepackage[es-nodecimaldot]{babel} % nombre de opción corregido
\usepackage{amsthm, amssymb, amsmath, amsfonts, bbm, mathtools}
\usepackage{csquotes}
\usepackage{enumerate}
\usepackage[shortlabels]{enumitem}
\usepackage{titling}
\usepackage{blindtext}
\usepackage{lmodern} 
\usepackage{float}
\usepackage{mathrsfs}
\usepackage[capitalise]{cleveref}
\usepackage{caption} 
\usepackage{graphicx}
\usepackage{chngcntr}
\usepackage{enumitem}
\usepackage[ruled,vlined]{algorithm2e}

\counterwithin{figure}{section}
\numberwithin{equation}{section}
\usepackage[vmargin=1in, hmargin=0.5in]{geometry}  

\setlength{\topmargin}{-50pt}
\setlength{\headheight}{38.23112pt}
\setlength{\headsep}{24pt}

\makeatletter
\def\iftitle{\ifx\thetitle \empty \else }
\def\safetitle{\iftitle{\thetitle}\fi}
\def\headertitle{%
    \ifdefined\hw{\hw}\fi%
    \ifdefined\hw \iftitle{ -- }\fi%
    \safetitle%
}
\def\@maketitle{%
  \begin{center}%
Tarea 4 - Problemas inversos - 25/03/2026 \\
Muñoz Macías Christian Jaime 
    \vskip 1em%
    \hrule
  \end{center}%
  \par
  \vskip 1.5em%
}
\makeatother
\usepackage{amsmath}
\usepackage{tikz}

\newcommand\RADot{%
  \mathord{%
    \mspace{1mu}%
    \text{\radot}%
    \mspace{1mu}%
  }%
}

\newcommand\radot{%
    \tikz[line cap=round,x=1ex,y=1ex,line width=0.3pt]
    {\draw (0,1) |- (1,0) (0.55,0) arc(0:90:0.55); \fill (0.23,0.23) circle (0.05);}%
}

%% encabezado personalizado
\usepackage{fancyhdr}
%% estilo plain personalizado
\fancypagestyle{plain}{%
    \fancyhf{} % limpia encabezado y pie de página
    \renewcommand{\headrulewidth}{0pt}
    \setlength{\headheight}{0pt}
    \fancyfoot[R]{\thepage}
}
%% estilo para todas las páginas
\pagestyle{fancy}
\fancyhf{}
\fancyhead[C]{Maestría en Matemáticas Aplicadas}
\fancyfoot[R]{\thepage}

%%%%%%%%%%%%%%%%%%%%%%%%
%% DEFINICIONES ÚTILES %%
%%%%%%%%%%%%%%%%%%%%%%%%
%% definiciones matemáticas
\newcommand{\N}{\mathbb{N}}
\newcommand{\Z}{\mathbb{Z}}
\newcommand{\Ha}{\mathcal{H}}
\newcommand{\R}{\mathbb{R}}
\newcommand{\Q}{\mathbb{Q}}
\newcommand{\C}{\mathbb{C}}
\newcommand{\PP}{\mathbb{P}}

\newtheorem{proposition}{Proposición}
\newtheorem{theorem}{Teorema}
\newtheorem{lemma}{Lema}
\newtheorem{corollary}{Corolario}

\theoremstyle{definition}
\newtheorem{problem}{Problema}
\newtheorem{definition}{Definición}[section]
\newtheorem*{example}{Ejemplo}

\theoremstyle{remark}
\newtheorem*{observation}{Observación}
\newtheorem*{note}{Nota}

\newenvironment{solution}{\begin{proof}[Solución]}{\end{proof}}
\newenvironment{dem}{\begin{proof}[Demostración]}{\end{proof}}
\renewcommand{\qedsymbol}{$\blacksquare$}

%% comandos útiles
\newcommand{\eps}{\varepsilon}
\newcommand{\To}{\longrightarrow}
\newcommand{\Eu}{\varPhi}
\newcommand{\Mapsto}{\longmapsto}

%% delimitadores emparejados
\DeclarePairedDelimiter{\set}{\{}{\}}
\DeclarePairedDelimiter{\floor}\lfloor\rfloor
\DeclarePairedDelimiter{\ceil}\lceil\rceil
\DeclarePairedDelimiter{\abs}\lvert\rvert
\DeclarePairedDelimiter{\gen}\langle\rangle
\DeclareMathOperator{\sgn}{sgn}
\DeclareMathOperator{\ord}{ord}

%% operador
\newcommand{\op}[1]{\operatorname{#1}}

%% etiquetas del documento
\title{Tarea 4 - Problemas inversos - 25/03/2026}
\def\course{Materia}
\def\university{CIMAT - MsC. Maestría en matemáticas aplicadas}
\author{Muñoz Macías Christian Jaime}
\def\email{\url{christian.munoz@cimat.mx}}
\date{\today}

%% permitir saltos de línea en ecuaciones
\allowdisplaybreaks[4]

\begin{document}

\maketitle

\section*{Ecuación de Poisson}
Consideremos el problema de valores a la frontera para la ecuación de Poisson
\[
-u''(x) = f(x), \quad x \in (0, L), \quad u(0) = u(L) = 0. \tag{1}
\]
La solución se puede escribir como una integral usando la función de Green $G(x, \xi)$:
\[
u(x) = \int_{0}^{L} G(x, \xi) f(\xi) \, d\xi,
\]
donde la función de Green está dada por:
\[
G(x, \xi) = \begin{cases} 
\dfrac{x(L - \xi)}{L}, & \text{si } x \le \xi, \\[10pt]
\dfrac{\xi(L - x)}{L}, & \text{si } x > \xi. 
\end{cases}
\]
Entonces, la solución general del problema es:
\[
u(x) = \int_{0}^{x} \frac{\xi(L-x)}{L} f(\xi) \, d\xi + \int_{x}^{L} \frac{x(L-\xi)}{L} f(\xi) \, d\xi. \tag{2}
\]

\section*{Ejercicios}

\begin{enumerate}
\item Implementa la cuadratura de la regla del punto medio compuesta para discretizar el problema (2) suponiendo que $L = 1$.
\begin{solution}
    definimos una $h$ la cual es 
    \[
    h=\frac{b-a}{n}
    \]
    donde n son las particiones del espacio, la cuadratura de la regla del punto medio compuesta esta definida como
    \[
    x_i^*=a+\left(i-\frac{1}{2}\right)h
    \]
    donde $x_i^*$ son los puntos de evaluación y mi integral es
    \[
    \int_{a}^{b}f(x)dx \approx h\sum_{i=1}^{n}f(x_i^*)
    \]
    entonces mi descretización de mi ecuación es

    la evaluacion de $\xi$ 
    \[
    \xi_i^*=0+\left(i-\frac{1}{2}\right)h
    \]
    
    donde $n_1$ y $n_2$ son $\mattbb{N}$ donde $n_1$ son los nodos del dominio de $\left[0,x\right]$ y $n_2$ los nodos del dominio $\left[x,L\right]$
    donde $u$ es
    \[
    u(x) \approx \sum_{i=1}^{n_1} \frac{\xi_{i}^{*}(L-x)}{L} f(\xi_{i}^{*}) h + \sum_{i=n_1}^{n_2} \frac{x(L-\xi_{i}^{*})}{L} f(\xi_{i}^{*}) h
    \]
    para problema especifico donde $L=1$

    \[
    u(x) \approx \sum_{i=1}^{n_1} \xi_{i}^{*}(1-x) f(\xi_{i}^{*}) h + \sum_{i=n_1}^{n_2} x(1-\xi_{i}^{*}) f(\xi_{i}^{*}) h
    \]
    donde se considera que $h$ es uniforme, donde $h=1/(n+1)$ 
    ahora si consideramos que se discretizar tambien respecto de x, entonces vamos tener dato de $u$, respecto a cada $j$-esimo 
    valor de $x$, vamos tener un vector $\textbf{x}$  y de $\textbf{u}$ y un vector de $\textbf{f}$
    esto se puede ver como sistema lineal como
    \[
    \textbf{u}(\textbf{x}) \approx A \textbf{f}
    \]  
    donde $A$ es
    \[
    \left[A\right]_{ij}
    =\begin{cases}
      \xi_{i}^{*}(1-x_j)h \quad & x_j \leq \xi_i \\
      x_j(1-\xi_{i}^{*}) h \quad & x_j \geq \xi_i
    \end{cases}
    \]
    donde $\textbf{f}$ es
    \[
    \left[\textbf{f}\right]_{i}=f(\xi_i)
    \]
    donde $\textbf{f}$ es mi incógnita 
\end{solution}
\item Haz un modelo a priori de $f(x)$ mediante un Gaussian Markov Random Field usando el artículo de Bardsley como referencia. Considera que la condición de frontera de $f$ es conocida. Nota: Al construir la distribución a priori, es necesario añadir un offset (desplazamiento) a la media; de lo contrario, la media a priori será cero y el modelo no capturará adecuadamente la información disponible sobre $f$.
\begin{solution}
  Se puede considerar que nuestro es un problema inverso lineal 
  \[
  b=Ax
  \]
  donde $b \in  \mathbb{R}^m$ que son los datos observados, donde $x$ es un vector $n \time 1$ vector desconocido y $A$ es una matriz de $m \time n$ es modelo . en la practica tenemos ruido.
  \[
  b=Ax+\eta
  \]
  donde $\eta$ es un vector random gaussiano. mi problema  es  una ecuacion Poisson
  \[
  b=Poiss(Ax+\gamma)
  \]
  donde $gamma$ es un vector $m \time 1$ de recuentos de fondo y se supone conocido.

  \[
  p(b|x) \propto \exp \left(-\frac{\lambda}{2}\|Ax-b\|^2\right)
  \]
  su forma logaritmica es 

  \[
  \log p(b|x) \propto  -\frac{\lambda}{2}\|Ax-b\|^2
  \]
  la probabilidad posteriori funcion densidad $p(x|b)$ se puede escribir como
  \[
  p(x|b) \propto p(b|x)p(x)
  \]
  su forma logaritmica es 

  \[
  \log p(x|b) \propto  \log p(b|x)+ \log p(x)
  \]
  el modelo a priori que proponen en el artículo de bardsley es
  \[
  P(x) \propto \exp\left(-\frac{\delta}{2}x^T L x\right)
  \]
  pero para este trabajo vamos a añadir una media $\mu$
  \[
  P(x) \propto \exp\left(-\frac{\delta}{2}(x-\mu)^T L (x-\mu)\right)
  \]
  donde $\delta L$ es matriz covarianza inversa. donde vector $\textbf{x}=(x_1,\dots, x_n)$ y \\
  donde $\textbf{x}_{-i}=(x_1,\dots,x_{i-1},x_{i+1},\dots,x_n)$
  \[
   \textbf{x}|\textbf{x}_{-i} \sim N \left(\sum_{j \neq i}\beta_{ij}x_j, \kappa_i^{-1} \right)
  \]
  los parametros sastifacen $\beta_{ij} \kappa_i=\beta_{ji} \kappa_j$ para todo $i$ y $j$
  \[
  p(x)=(2 \pi)^{-n/2}|Q|^{1/2}\exp \left(-\frac{1}{2}(x-\mu)^T Q (x-\mu)\right)
  \]
  su forma logaritmica es
  \[
  \log p(x)=-\frac{n}{2} \log(2 \pi)+ \frac{1}{2}\log|Q| -\frac{1}{2}(x-\mu)^T Q (x-\mu)
  \]
  donde $|Q|$ es el $\det(Q)$ es
  \[
  \det|Q|=\prod_{i=1}^{n} \lambda_i
  \]
  si calculamos su logaritmo
  \[
  \log \det|Q|=\sum_{i=1}^{n} \log \lambda_i
  \]
  donde 
  \[
  \left[Q \right]_{ij}=\begin{cases}
    \kappa_i,  \quad & i=j \\
    -\kappa_i \beta_{ij} \quad & i \neq j
  \end{cases}
  \]
  La malla computacional que se utiliza para realizar la discretización numérica define un conjunto de nodos i y un 
  sistema de vecindad  $\partial_i$ para $i = 1,\dots,n$. Resulta que los elementos de $\partial_i$ son los nodos $j = i$ para los cuales $\left[Q\right]_{ij} = 0$, 
  y por lo tanto,
  \[
  \partial_i  =\{j \neq i| \beta_{ij} \neq 0\}
  \]
  ademas, si $\textbf{x}_{\partial_i}=\{x_j | j \in \partial_i \}$, entonces $p\left(x_i|x_{-i}\right)=p\left(x_i|x_{\partial_i}\right)$
  Como ejemplo, supongamos que
\[
\kappa_i = \frac{\delta n_i}{h^2},
\]
donde \( n_i = |\partial i| \) y \( \delta > 0 \), y
\[
\beta_{ij} =
\begin{cases}
\frac{1}{n_i}, & \text{si } j \in \partial i, \\
0, & \text{en otro caso}.
\end{cases}
\]

Entonces, toma la forma
\[
x_i \mid x_{\partial i}
\sim \mathcal{N}\left(\bar{x}_{\partial i}, \left(\frac{\delta n_i}{h^2}\right)^{-1}\right),
\]
donde
\[
\bar{x}_{\partial i} = \frac{1}{n_i} \sum_{j \in \partial i} x_j
\]
es el promedio de los valores vecinos, y la función de densidad conjunta está dada por
\[
\left[Q \right]_{ij} = \frac{\delta}{h^2}
\begin{cases}
n_i, & i = j, \\
-1, & j \in \partial i, \\
0, & \text{en otro caso}.
\end{cases}
\]
\end{solution}
\item Implementa la solución bayesiana del problema inverso de determinar
\[
f(x) = x
\]
dadas observaciones ruidosas de
\[
u(x) = \frac{1}{6}x(1-x)(1+x)
\]
con ruido aditivo gaussiano con relación de señal a ruido 100.
\begin{solution}

\end{solution}
\item Calcula la media y la covarianza posterior de forma analítica.

\item Muestra tu distribución posterior y compara la media posterior numérica tomando en cuenta el tiempo integrado de autocorrelación.

\item Calcula la constante de normalización de manera analítica, y numéricamente de la manera siguiente:

\begin{enumerate}
\item[(a)] Usa la función \texttt{scipy.stats.gaussian\_kde} para ajustar una distribución gaussiana $G(f)$ a las muestras que obtuviste.

\item[(b)] Aproxima
\[
\begin{aligned}
z_0 &= \int \pi(u|f)\pi(f)\,df \\
&= \int \frac{\pi(u|f)\pi(f)}{G(f)} G(f)\,df \\
&\approx \frac{1}{N} \sum_{j=1}^{N} \frac{\pi(u|f_j)\pi(f_j)}{G(f_j)}
\end{aligned}
\]
donde $\{f_j\}_{j=1}^N$ son muestras de $G(f)$.
\end{enumerate}
\end{enumerate}

\section*{El Truco del Log-Sum-Exp: Una Herramienta para la Estabilidad Numérica}

\subsection*{Motivación}
En muchos cálculos probabilísticos, frecuentemente nos encontramos con expresiones del tipo:
\[
\log \left( \sum_{i=1}^{N} e^{a_i} \right)
\]
Esta expresión aparece naturalmente al calcular el logaritmo de una mezcla, log-verosimilitud de modelos marginales, y en particular al estimar esperanzas en espacio logarítmico como:
\[
\log \left( \frac{1}{N} \sum_{i=1}^{N} w_i \right) = \log \left( \frac{1}{N} \sum_{i=1}^{N} e^{\log w_i} \right)
\]
El cálculo directo de esta cantidad puede sufrir de inestabilidad numérica debido a desbordamientos (overflow) o subdesbordamientos (underflow) cuando los exponentes $a_i$ son muy positivos o muy negativos.

\subsection*{El Truco}
Sean $a_1, a_2, \ldots, a_N \in \mathbb{R}$. Definimos:
\[
M = \max_i a_i
\]
Entonces:
\[
\log\left(\sum_{i=1}^{N} e^{a_i}\right) = M + \log\left(\sum_{i=1}^{N} e^{a_i - M}\right)
\]

\subsection*{¿Por qué Funciona?}
Dado que $a_i - M \leq 0$, los exponenciales $e^{a_i - M}$ son todos menores o iguales a 1, evitando desbordamientos. Además, si los $a_i$ son muy negativos, restar $M$ lleva el mayor término a cero, evitando subdesbordamientos en los otros términos.

\subsection*{Ejemplo}
Supongamos:
\[
a = [-1000, -1001, -999]
\]
Entonces:
\[
\sum_{i=1}^{3} e^{a_i} = e^{-1000} + e^{-1001} + e^{-999}
\]
Estos números son demasiado pequeños para ser representados en punto flotante estándar. Usando el truco:
\[
\log\left(\sum_{i=1}^{3} e^{a_i}\right) = -999 + \log\left(e^{-1} + e^{-2} + 1\right) \approx -999 + \log\left(1 + e^{-1} + e^{-2}\right)
\]

\subsection*{En la Práctica}
En Python (por ejemplo, usando NumPy), esto se implementa como:
\begin{verbatim}
def logsumexp(a):
    a_max = np.max(a)
    return a_max + np.log(np.sum(np.exp(a - a_max)))
\end{verbatim}
Muchas bibliotecas científicas como \texttt{scipy} y \texttt{PyTorch} también implementan una función estable llamada \texttt{logsumexp}.

\end{document}